\documentclass[conference]{IEEEtran}
\usepackage[utf8]{inputenc}
\usepackage[T1]{fontenc}
\usepackage{cite}
\usepackage{amsmath,amssymb}
\usepackage{graphicx}
\usepackage{booktabs}
\usepackage{url}

\title{Verifier‑Assisted versus One‑Shot LLM Intent\ensuremath{\rightarrow}Configuration: A Paired NetOpsTraceBench Study}

\author{
\IEEEauthorblockN{Autonomous AI Researcher}
\IEEEauthorblockA{CNSM 2026 Agentic AI Researcher Track}
}

\begin{document}

\maketitle

\begin{abstract}
We preregistered and executed NetOpsTraceBench, a paired comparison of a one‑shot LLM intent\ensuremath{\rightarrow}configuration pipeline and a verifier‑assisted generate--validate--repair pipeline on 1,200 intent\ensuremath{\rightarrow}config episodes. Using gpt-5-mini, a deterministic NetOps validator, and Mininet‑style emulation, we applied a preregistered binary episode success rule (validator accepts final config AND emulator shows no availability degradation above tolerance AND no automatic rollback). Of 1,200 planned paired tasks we completed 1,199 pairs. Baseline success = 1,196/1,199 (0.9974979149); guarded success = 1,199/1,199 (1.0000000000). The paired‑success‑rate difference (guarded - baseline) = 0.0025020851 (95\% bootstrap percentile CI [0.0000, 0.0058381985], bootstrap\_seed=1729, resamples=10,000); exact McNemar two‑sided p = 0.25. The preregistered power target sought a 5 percentage‑point minimum detectable effect; given the observed near‑ceiling baseline, the study lacked power to detect much smaller absolute gains. Representative validator traces and provider audit logs show repair was invoked only three times. Under these controlled emulation and task‑corpus conditions we find no statistically significant advantage for the verifier‑assisted pipeline; we report artifacts, analytic deviation (Wilcoxon \ensuremath{\rightarrow} exact McNemar for binary outcomes), and detailed execution accounting to support reproducible interpretation.
\end{abstract}

\section{Introduction}
Generative large language models (LLMs) are increasingly proposed to translate operator intent expressed in natural language into executable device configurations for network operations (NetOps). Systems and survey literature emphasize that operational reliability depends critically on orchestration, deterministic validators, and assurance infrastructure rather than on an LLM alone [1,2]. Generate--validate--repair (verifier‑assisted) architectures are a canonical mitigation strategy: an LLM proposes candidate outputs, an automated verifier checks constraints and safety, and an automated repair loop iterates until the candidate satisfies verifier rules [3--5]. Despite conceptual support, publicly archived, preregistered, paired evidence measuring whether verifier assistance yields an operationally meaningful improvement (rollback‑aware binary success) relative to one‑shot generation is sparse.

We address this gap with NetOpsTraceBench: a preregistered, paired, emulator‑grounded study comparing (A) baseline one‑shot NL\ensuremath{\rightarrow}config generation and (B) a guarded generate--validate--repair pipeline using the same shared initial candidate. The study asked: Can a reproducible, workflow‑centered benchmark of paired intent\ensuremath{\rightarrow}config episodes produce an objective paired binary success metric that reliably distinguishes verifier‑assisted pipelines from one‑shot generation across two simulated topologies? Our contributions are (1) a preregistered experimental design and full execution artifacts (study ID rX4f7-1); (2) a paired analysis on 1,199 completed pairs using gpt-5-mini and a deterministic validator; (3) transparent reporting of an analytic deviation (final confirmatory test used an exact McNemar test appropriate for paired binary outcomes) and execution accounting that explains why the guarded pipeline rarely required repair.

\section{Related Work}
Recent survey and system literature frames LLMs as components inside broader AIOps/NetOps workflows and repeatedly emphasizes that correctness and trustworthiness are properties of the whole orchestration stack (model + verifier + provenance + sandboxed execution), not of the base LLM alone [1--3]. Multiple community reviews urge workflow‑centered evaluation---trace quality, validator interactions, and rollback‑aware metrics---rather than isolated one‑shot QA tests; these reviews motivate reproducible, audit‑ready benchmarks that combine generation, deterministic validation, and emulation [1,2,4].

Concrete verifier‑assisted approaches appear across adjacent domains and NetOps prototypes. Work on NL\ensuremath{\rightarrow}formal translations and CTL/specification generation demonstrates generator+verifier loops where formal checks drive iterative repairs; similarly, topology‑compilation systems translate high‑level security intent into executable topologies with compliance checks and constraint enforcement [5,6]. Prototype efforts in the configuration space (text\ensuremath{\rightarrow}config generators and misconfiguration detectors) and domain‑specific "LLM firewalls" illustrate practical patterns: an LLM proposes a candidate, an automated validator enforces syntactic and semantic constraints, and an orchestrator optionally triggers constrained repair or rejects the proposal [7--10]. These systems commonly instrument deterministic parsers and domain rules to produce binary verifier outcomes (accept/reject) and produce their own evaluation sets for demonstration; as a result, reported gains are often conditional on task selection and validator strictness.

Empirical and conceptual studies document LLM failure modes that motivate guarded architectures. Recent analyses identify hallucination, unit/semantic errors, overconfidence/miscalibration, provenance loss, and susceptibility to prompt‑injection as recurrent risks when LLMs propose executable artifacts---risks that an external verifier or provenance linker can detect or mitigate in some instances but not eliminate universally [11--14]. Several papers demonstrate that verifier‑guided repair reduces syntactic and some semantic errors in constrained settings, but they also highlight sensitivity to the verifier's rule set (strictness) and to the representativeness of evaluation examples. In short, prior work shows feasibility of generate--validate‑repair patterns but also underscores the dependence of measured improvement on corpus difficulty, adversarial content, and verifier configuration.

Related benchmarking and dataset discussions emphasize gaps that NetOpsTraceBench targets. System papers often construct private or narrowly scoped held‑out sets calibrated to their prototype; TopoIntent and other recent tools note the absence of large public, vendor‑diverse intent\ensuremath{\rightarrow}config datasets with execution ground truth and rollback semantics, which complicates cross‑system comparisons [6,7,9]. Surveys therefore call for standardized, workflow‑centered benchmarks that (a) record full provider and validator traces, (b) embed emulator execution to check operational side effects (availability/rollback), and (c) preregister analysis choices to reduce researcher degrees of freedom [1,2]. NetOpsTraceBench operationalizes these recommendations by preregistering inclusion/exclusion rules, seeds, a binary, rollback‑aware episode success metric, and by archiving raw model outputs, validator traces, and execution manifests so that the effect of verifier design and task corpus can be examined reproducibly.

Taken together, the literature supports the conceptual value of verifier scaffolding while also documenting two recurring caveats that our study made explicit in design and interpretation: (i) measured gains depend strongly on task difficulty and verifier thresholds, and (ii) many demonstrations use constructed datasets and domain‑specific validators, which limits generality. NetOpsTraceBench contributes a preregistered, artifact‑rich paired evaluation that makes these dependencies observable (repair invocation counts, contingency discordances, and emulator rollback checks) rather than inferred, enabling follow‑up sensitivity studies that vary corpus adversariality and verifier strictness.

\section{Methodology}
Preregistration and experimental scope: The study design was preregistered under ID rX4f7-1 and frozen prior to execution (preregistration manifest SHA‑256 = 77abe8b4...). The preregistration specified confirmatory hypotheses, a single primary estimand (paired\_success\_rate\_difference\_guarded\_minus\_baseline), deterministic inclusion/exclusion rules, contamination thresholds, stopping rules, random seeds, maximum model‑call budget, and a complete analysis plan. NetOpsTraceBench\_v1 targeted 1,200 paired intent\ensuremath{\rightarrow}config tasks partitioned into two simulated topology clusters (edge = 600, core = 600). Task examples were public synthetic or consented; every task carried deterministic provenance metadata. The preregistered power plan assumed a baseline success near 0.70 and set a minimum detectable effect (MDE) of 5 percentage points (absolute) at two‑sided \ensuremath{\alpha}=0.05 and power 0.8; sensitivity scenarios over a range of baseline and discordant‑pair assumptions were specified in the preregistration.

Pipelines, transformations, and instrumented components: We compared two paired conditions executed from a single shared initial candidate per pair. Baseline (one‑shot) invoked a single LLM generation call to produce an initial candidate configuration. Guarded (verifier‑assisted) used the same shared initial candidate and permitted at most one automated repair iteration when the deterministic validator rejected the candidate; preregistration limited maximum\_repair\_calls\_per\_task = 1 and initial\_generation\_calls\_per\_task = 1 to bound model‑call budget. The requested model was gpt-5-mini (provider recorded as openai\_responses) invoked through orchestration adapter hosted\_netops\_gvr\_v1. Deterministic\_netops\_task\_generator\_v5 produced canonical, vendor‑agnostic intent\ensuremath{\rightarrow}configuration tasks; deterministic\_netops\_validator\_v5 performed parsing, intent‑constraint validation, dependency checks, and emulator preflight evaluation. The execution contract capped maximum\_model\_calls = 2,400 and defined the scientific\_confirmatory execution mode.

Inclusion, contamination, and stopping rules: Inclusion required that neither episode in a pair be deterministically excluded. The preregistration implemented contamination detection by (a) a normalized substring/token‑overlap test with threshold T\_contam = 0.85 and (b) an optional embedding‑similarity secondary check (cosine threshold = 0.92) for flagged items. If either test exceeded thresholds the pair was flagged and excluded from the primary confirmatory set (but retained for exploratory reporting). Pairs were also excluded if both episodes failed to produce parsable vendor renderings after one automated re‑execution attempt. The stopping rule required running the full planned workload unless (i) cumulative model calls reached the maximum\_model\_calls cap, or (ii) automation halted due to deterministic component failure after the one allowed automated recovery attempt; all such events were logged.

Determinism, seeds, and failure handling: The preregistration froze generation seeds and analysis seeds. Deterministic task generation and validator behavior reduce sampling variance from non‑deterministic test artifacts; model nondeterminism remained possible at the provider side but was bounded by the single‑call-per‑stage policy. Missingness handling followed the preregistered plan: the primary analysis used complete\_pair\_primary (exclude flagged/excluded pairs); a sensitivity analysis treated excluded or failed episodes as failures (complete\_pair\_primary\_with\_failure\_as\_zero\_sensitivity). All exclusion reasons, failed episode counts, and provider audit traces were archived to preserve transparent accounting.

Success metric, confirmatory test, and bootstrap CI: The preregistered binary episode success rule required three conjunctive outcomes: (1) final configuration accepted by the deterministic validator (parsing and intent constraints satisfied), (2) emulator execution shows availability degradation no greater than the preregistered tolerance, and (3) no automatic rollback event occurred during simulated execution. Although the preregistration specified a paired nonparametric Wilcoxon signed‑rank test, the archived outcomes were binary for each episode. Following the preregistered emphasis on correct paired‑outcome inference, the final archived confirmatory analysis applied an exact McNemar two‑sided test to the 2\ensuremath{\times}2 paired contingency table and computed a paired bootstrap percentile 95\% confidence interval for the paired‑success‑rate difference using frozen bootstrap\_seed = 1729 and 10,000 resamples. This analytic deviation (Wilcoxon \ensuremath{\rightarrow} exact McNemar) is recorded in the analysis artifacts and in the Disclosure Statement.

Instrumentation, auditing, and artifact recording: The run captured provider call audit metadata (hosted\_model\_call\_event\_count and stage\_counts), all raw model outputs, shared initial candidate traces, repair provider traces when present, validator traces (validation\_before/validation\_after), emulator execution logs, task\_manifest entries, transformation\_manifest, and a full hash manifest. Key configured limits recorded in model\_configuration.json included max\_output\_tokens = 1500 and requested\_model = gpt-5-mini. The execution manifest records planned\_episode\_count = 2,400, completed\_episode\_count = 2,398, failed\_episode\_count = 2, and model\_calls\_used = 1,203. All seeds, analysis code (paired\_binary\_analysis\_v1), and final analysis artifacts (analysis/paired\_contingency\_table.csv, analysis/results.json, analysis/paired\_difference\_figure.svg) are archived to enable exact reproduction.

Predefined sensitivity and exploratory plans: The preregistration enumerated exploratory analyses that the archived run executed or preserved for downstream work: (a) distribution of repair iterations and their correlation with final success and time‑to‑repair; (b) automated failure‑mode categorization (semantic unit errors, missing dependencies) derived from deterministic parsers and emulator checks; (c) sensitivity of the primary binary success to alternative availability‑degradation tolerances and to treating flagged contamination pairs as failures. All exploratory outputs were labeled as such and are reported separately from the confirmatory inference to control Type I error per the preregistered multiplicity plan.

\section{Results}
Execution and primary accounting (archived artifacts). The run attempted the preregistered 2,400 episodes (1,200 paired tasks) and completed 2,398 episodes (completed\_pairs = 1,199; failed\_episode\_count = 2). Provider audit fields recorded hosted\_model\_call\_event\_count = 1,203, completed\_hosted\_model\_call\_event\_count = 1,202, and failed\_hosted\_model\_call\_event\_count = 1. Stage-level accounting shows shared\_initial\_generation = 1,200 and repair = 3 (execution/provider audit and deterministic\_reconciliation.json). Artifact\_count for the full run = 8,408 and the run completed at 2026-08-30T03:15:12.141890+00:00 (execution manifest).

Primary confirmatory outcomes (archived analysis outputs). Using the preregistered binary success rule, the archived condition summary (analysis/condition\_summary.csv) reports baseline\_success\_count = 1,196 of 1,199 complete pairs (baseline\_success\_rate = 0.9974979149291076) and guarded\_success\_count = 1,199 of 1,199 (guarded\_success\_rate = 1.0). The paired contingency table (analysis/paired\_contingency\_table.csv) is n11 = 1,196 (both succeeded), n10 = 3 (guarded succeeded, baseline failed), n01 = 0 (baseline succeeded, guarded failed), n00 = 0 (both failed). The point estimate for the paired success‑rate difference (guarded - baseline) = 0.002502085070892446. The archived confirmatory exact McNemar two‑sided p-value = 0.25. The paired bootstrap percentile 95\% confidence interval computed with bootstrap\_seed = 1729 and 10,000 resamples is [0.0000, 0.005838198498748957] (analysis/results.json and analysis/analysis\_log.jsonl).

Missingness and failed-call accounting. The preregistered missingness summary and deterministic reconciliation artifacts report complete\_pair\_count = 1,199, failed\_baseline\_episodes = 1, failed\_guarded\_episodes = 1, and both\_missing\_pairs = 1 (analysis/missingness\_summary.csv; deterministic\_reconciliation.json). No pairs were excluded for contamination under the preregistered thresholds (analysis/contamination\_summary.csv empty). The run therefore uses the complete\_pair\_primary treatment as planned.

Repair invocation and provider audit diagnostics. Repair-stage calls were invoked only three times (stage\_counts.repair = 3 in provider\_call\_audit). Provider audit also shows no cross-condition cache reuse (cross\_condition\_cache\_reuse\_observed = false) and confirms the intended shared\_initial\_generation semantics (shared initial candidate used for paired episodes). The low repair count is corroborated by the scoring artifacts: representative scoring JSONs (e.g., execution/scoring/task-000001-baseline.json and execution/scoring/task-000001-guarded.json) show validator\_trace.validation\_after.valid = true and validator\_trace.repair\_applied = false for typical episodes. In other words, the deterministic validator accepted the initial shared candidate in the vast majority of tasks, leaving only three guarded episodes that required a repair call.

Representative per-task diagnostics (artifact-grounded). Representative task\_manifest entries and scoring traces illustrate typical episode structure and validator acceptance criteria. For example, pair-000001 (execution/scoring/task-000001-baseline.json) had parsed\_command\_count = 4, required\_command\_count = 4, observed\_touched\_field\_count = 3, and validation\_after.normalized\_configuration matching the reference; validation\_after.valid = true and repair\_applied = false. A harder task example (pair-000002) shows difficulty.level = "hard" with changed\_interface\_count = 2 and required\_command\_count = 2, yet its initial candidate was still validated as valid by the deterministic validator. These per-episode fields (parsed\_command\_count, required\_command\_count, difficulty pattern) are recorded in the representative scoring artifacts and demonstrate that many tasks contained explicit required sequences which the generator reproduced correctly.

Interpretation of discordant pairs and CI. The observed discordant pattern (n10 = 3, n01 = 0) implies that all discordant evidence favored the guarded pipeline; however, the absolute number of discordant pairs is tiny relative to the total sample. The bootstrap CI lower bound is exactly 0.0000 (percentile method) and the upper bound \ensuremath{\approx} 0.00584, indicating that the data are compatible with a small positive absolute gain but not with the preregistered target MDE (5 percentage points). The exact McNemar p = 0.25 formalizes that, under these discordant counts, the test does not reject the null of no paired difference at \ensuremath{\alpha} = 0.05.

Why repair was rarely exercised. Two archived facts explain the infrequent repair calls: (1) the synthetic task generator produced tasks with explicit required sequences and constrained allowed field changes (deterministic\_netops\_task\_generator\_v5 task\_payloads record required\_sequence entries), increasing the chance that a straightforward canonical candidate satisfies intent constraints; (2) the deterministic validator configuration (deterministic\_netops\_validator\_v5) accepted initial candidates when all intent constraints and preflight checks matched the rendered candidate. Provider audit stage\_counts (shared\_initial\_generation = 1,200, repair = 3) and scoring artifacts jointly document these operational mechanisms and support the conclusion that the guarded pipeline had little opportunity to improve over an already‑high baseline under the chosen corpus and validator strictness.

Ancillary archived metrics. The analysis artifact set includes analysis/paired\_difference\_figure.svg (visual CI and point estimate), analysis/paired\_contingency\_table.csv (full contingency counts), and analysis/analysis\_log.jsonl (bootstrap parameters: resamples = 10,000; seed = 1729). The execution manifest archives timing, provider call counts, and the preregistration SHA‑256 reference; these artifacts are sufficient to reproduce the confirmatory estimates and the diagnostic counts reported above.

\section{Discussion}
The statistical evidence and operational diagnostics together explain why the preregistered paired comparison did not find a statistically significant advantage for the guarded pipeline in this run. The exact McNemar two‑sided test (p = 0.25) and the paired bootstrap percentile 95\% CI [0.0000, 0.0058381985] both include zero, indicating that the observed paired success‑rate difference (\ensuremath{\approx}0.25 percentage points) is consistent with sampling variability under the recorded discordant counts (n10 = 3, n01 = 0). Those discordant counts are the proximal determinant of inference here: with 1,199 complete pairs, only three pairs were discordant and all favored the guarded pipeline, which yields a point estimate but insufficient evidence to reject the null at \ensuremath{\alpha}=0.05. This arithmetic --- small discordant counts in an otherwise near‑ceiling outcome distribution --- is why the preregistered MDE of 5 percentage points (chosen under an assumed baseline \ensuremath{\approx}0.70) was not matched by the observed data; the study therefore lacked confirmatory power to detect such a small absolute improvement.

Operationally, archived provider and validator artifacts clarify the mechanism that produced these statistics. The provider audit reports model\_calls\_used = 1,203 with stage\_counts showing shared\_initial\_generation = 1,200 and repair = 3, and representative validator traces (e.g., execution/scoring/task-000001-baseline.json) show validation\_after.valid = true for typical episodes. Together these elements indicate that the shared initial candidates satisfied the deterministic\_netops\_validator\_v5 checks in the overwhelming majority of episodes and so the guarded pipeline's repair stage was rarely invoked. Thus, under the exact combination of (a) a synthetic, well‑specified task generator that produces tasks with explicit required sequences, and (b) a validator configuration that enforces those sequences in a way that treats correct initial candidates as acceptable, the marginal operational value of generate--validate--repair for the binary success metric is necessarily small.

These findings do not falsify the general utility of verifier‑assisted architectures. Instead they highlight two contextual dependencies that practitioners and benchmark designers must consider. First, baseline difficulty and task ambiguity directly determine headroom for improvement: when initial generation already meets deterministic verifier criteria, measured gains on a strict binary success metric will be limited. Second, verifier strictness and the mapping between validator checks and operational safety matter: different validator thresholds or richer emulation of transient behaviors (e.g., routing convergence, vendor‑specific command ordering) would change repair invocation rates and the distribution of discordant pairs. Both dependencies are visible in the archived run artifacts (task\_manifest entries and scoring traces) and can be explored by varying task ambiguity, contamination thresholds, or validator rules using the provided analysis code and seeds.

Practical implications follow for benchmark design and operational evaluation. To measure verifier benefit robustly, a benchmark should (i) include a larger proportion of ambiguous or adversarial intents that create meaningful syntactic/semantic failure modes for LLMs, (ii) vary verifier strictness (from permissive to conservative) to expose where repairs would change execution outcomes, and (iii) exercise multiple model families and rendering targets to assess generality across vendor syntaxes. The preregistered sensitivity scenarios and the archived artifacts make such follow‑up experiments reproducible: the run archives deterministic seeds, task generator configuration, validator traces, and analysis scripts so other researchers can replay alternative verifier sensitivities or inject adversarial perturbations without ambiguity.

Cost--benefit considerations also emerge from the execution accounting. Repair was invoked only three times, and the total model call audit (\ensuremath{\approx}1,203 calls) shows a small additional invocation footprint relative to the number of pairs. In settings where repairs are rare because validators accept correct initial candidates, the operational overhead of a guarded pipeline may be modest; however, this mechanical accounting does not capture human‑in‑the‑loop costs, potential latency requirements, or the upstream effort to build and maintain robust validators and emulators. Those broader operational costs and benefits remain important practical considerations beyond the binary success metric evaluated here.

Finally, the run exposes reproducible levers for future investigation. The archived CI, contingency table, and provider logs document both the null‑result boundary conditions and the exact artifacts required to test alternative hypotheses (e.g., higher task ambiguity, stricter validators, adversarial stress generators, multi‑model comparisons). Because the preregistration and execution artifacts are public within this evidence bundle, researchers can systematically vary one factor at a time and quantify how discordant counts and the paired‑difference distribution change --- an approach that will be necessary to move from context‑specific null findings toward generalizable operational guidance about when verifier scaffolding materially improves NetOps safety.

\section{Limitations}
\begin{itemize}
\item Single model family tested (gpt-5-mini); results do not imply generality across other LLMs or fine‑tuned variants.
\item Task corpus skew: synthetic/consented tasks yielded near‑ceiling baseline success (\ensuremath{\approx}99.75\%), limiting headroom and statistical power to detect practical improvements by guarded pipelines.
\item Emulator fidelity: Mininet‑style emulation and deterministic validators approximate production behavior but may not capture all deployment failure modes (routing convergence, vendor idiosyncrasies, transient telemetry).
\item Verifier configuration dependence: deterministic\_netops\_validator\_v5 acceptance thresholds and parsing rules materially affect outcomes; different verifier strictness could change repair invocation rates and measured gains.
\item Analytic deviation documented: the preregistration specified a Wilcoxon signed‑rank paired procedure; the archived confirmatory analysis used an exact McNemar test because outcomes were binary. This deviation is recorded in the archived analysis artifacts and in the Disclosure Statement above.
\item Operational external validity: absence of heterogeneous vendor renderers, real production templates, and live rollouts constrain direct production generalization.
\end{itemize}

\section{Conclusion}
We preregistered and executed NetOpsTraceBench, a paired, emulator‑grounded study comparing one‑shot and verifier‑assisted NL\ensuremath{\rightarrow}configuration pipelines on 1,199 completed pairs. Under our synthetic task corpus, deterministic validator settings, and gpt-5-mini as requested, baseline one‑shot generation already achieved near‑ceiling success (\ensuremath{\approx}99.75\%). The guarded pipeline increased absolute success by \ensuremath{\approx}0.25 percentage points (paired difference = 0.0025020851) with exact McNemar two‑sided p = 0.25 and a 95\% bootstrap CI that includes zero. Archived execution manifests, provider audits (model\_calls\_used = 1,203; repairs invoked = 3), representative validator traces, and analysis artifacts support the interpretation that the observed null result reflects limited headroom given task and verifier choices rather than definitive evidence that verifiers are never useful. Future work should intentionally design more challenging and adversarial task sets, vary verifier strictness, and test multiple LLM families to determine conditions under which verifier scaffolding yields operationally meaningful improvements. All run artifacts, analysis code, seeds, and logs are archived to enable exact reproduction and follow‑up sensitivity studies (see Disclosure Statement).

\section*{Disclosure Statement}
Mandatory disclosure (preregistered run rX4f7-1). LLM(s) and provider: requested model = gpt-5-mini via provider recorded as openai\_responses; model configuration archived at execution/model\_configuration.json (requested\_model=gpt-5-mini, max\_output\_tokens=1500). Orchestration/adapter: hosted\_netops\_gvr\_v1 executed the paired benchmark. Deterministic tooling: deterministic\_netops\_validator\_v5 performed canonical parsing, intent‑constraint validation, and emulator preflight checks; deterministic\_netops\_task\_generator\_v5 produced synthetic tasks. Exact initial master prompt: archived at provenance/master\_prompt.txt (immutable reference SHA‑256 = 1872df1e\allowbreak{}1805d2d9\allowbreak{}6940456c\allowbreak{}a016bd66\allowbreak{}5d1d5196\allowbreak{}add77f5a\allowbreak{}cdf1582b\allowbreak{}b39b15ba). Topic constraint and autonomy boundary: the human Research Director selected the topic family and approved the master prompt pre‑lock; after the autonomy lock the AI autonomously formulated the research question, conducted literature review, generated the preregistration, executed the experiment, performed analysis, and wrote and revised the manuscript. Preregistration: NetOpsTraceBench (ID rX4f7-1) preregistration manifest SHA‑256 = 77abe8b4\allowbreak{}93bde444\allowbreak{}84e0b3d0\allowbreak{}2cfd7f9e\allowbreak{}c1621ff7\allowbreak{}e7847889\allowbreak{}082ae177\allowbreak{}a7ed5712; preregistered design (confirmatory hypotheses, inclusion/exclusion, seeds, stopping rules, analysis plan) was frozen before execution. Execution summary: started 2026-08-30T01:21:55.415955+00:00, completed 2026-08-30T03:15:12.141890+00:00; planned episodes 2,400; completed episodes 2,398; completed pairs 1,199; model\_calls\_used = 1,203; repair-stage calls observed = 3. Analysis artifacts and deviation record: the preregistration specified a Wilcoxon signed‑rank paired procedure; the archived executed confirmatory analysis used an exact McNemar two‑sided test on paired binary outcomes plus a paired bootstrap percentile CI. This post‑preregistration analytic choice is documented in the archived analysis artifacts (analysis/paired\_contingency\_table.csv, analysis/results.json, analysis/analysis\_log.jsonl, analysis/paired\_difference\_figure.svg) produced as part of the run that completed at 2026-08-30T03:15:12.141890+00:00. Artifacts and reproducibility: raw model outputs, provider traces, validator traces, task\_manifest, transformation\_manifest, execution logs, analysis code, and all seeds are archived; artifact\_count = 8,408 and full hash manifest path = execution/execution\_manifest.json. Human involvement: pre‑lock collaboration and infrastructure development occurred; no human manual editing, selective revision, or ad hoc text insertion occurred on the manuscript content after the autonomy lock---the AI performed internal autonomous revision cycles only. Technical restarts and deviations: execution manifest records two failed episodes and the analytic deviation described above; all deviations are logged in execution/execution\_log.jsonl. The disclosure above is complete relative to the archived artifacts supplied with this run.

\begin{thebibliography}{99}
\bibitem{ref1} http://hdl.handle.net/20.500.12708/229494
\bibitem{ref2} https://arxiv.org/abs/2608.13389
\bibitem{ref3} http://arxiv.org/abs/2403.09369
\bibitem{ref4} https://doi.org/10.2139/ssrn.6947162
\bibitem{ref5} https://doi.org/10.2139/ssrn.7222278
\bibitem{ref6} https://doi.org/10.20935/acadai8260
\bibitem{ref7} https://doi.org/10.3390/info17030297
\bibitem{ref8} https://doi.org/10.3390/app16010085
\end{thebibliography}

\end{document}
